\documentclass[../main]{subfiles}
\begin{document}
\chapter{Sistema trif\'asico{} desbalanceado}
\section{C\'alculo{} de corriente de neutro}
\subsection{Introducci\'on{}}
Un sistema trif\'asico{} es un tipo de sistema polif\'asico{} que utiliza tres se\~nales{} de corriente alterna que est\'an{} desfasadas entre s\'i{} $120^\circ$. Estos sistemas son fundamentales en la distribuci\'on{} de energ\'ia{} el\'ectrica{} y pueden clasificarse en sistemas balanceados y desbalanceados. En un \textbf{sistema trif\'asico{} balanceado}, todas las corrientes y tensiones son iguales en magnitud y est\'an{} desfasadas exactamente $120^\circ$ entre s\'i{}. Sin embargo, en un \textbf{sistema trif\'asico{} desbalanceado}, las magnitudes o los \'angulos{} de fase de las corrientes y tensiones var\'ian{}, lo que genera una corriente de neutro.

\info{Corriente de Neutro}{La corriente de neutro surge cuando las corrientes de las fases no se cancelan entre s\'i{} debido al desbalance. Esta corriente puede calcularse mediante la suma vectorial de las corrientes de las fases.}

\subsection{F\'ormulas{} fundamentales}
Para calcular la corriente de neutro en un sistema trif\'asico{} desbalanceado, utilizamos la siguiente expresi\'on{}:
\begin{equation}
I_N = \sqrt{I_A^2 + I_B^2 + I_C^2 - 2 \cdot I_A \cdot I_B \cdot \cos(120^\circ) - 2 \cdot I_B \cdot I_C \cdot \cos(120^\circ) - 2 \cdot I_C \cdot I_A \cdot \cos(120^\circ)}
\end{equation}
Donde:
\begin{itemize}
\item $I_N$: Corriente de neutro \unit{A}
\item $I_A, I_B, I_C$: Corrientes de las fases A, B y C respectivamente \unit{A}
\end{itemize}

\subsection{Ejemplo}
\ejemplo{Ejemplo}{Calcular la corriente de neutro en un sistema trif\'asico{} desbalanceado donde:
\begin{itemize}
\item $I_A = \qty{10}{A}$
\item $I_B = \qty{8}{A}$
\item $I_C = \qty{7}{A}$
\end{itemize}
Usamos la f\'ormula{} para la corriente de neutro (considerando desfase de $120^\circ$ y $\cos(120^\circ) = -0,5$):
\begin{align*}
I_N &= \sqrt{10^2 + 8^2 + 7^2 - 2 \cdot (10 \cdot 8 + 8 \cdot 7 + 7 \cdot 10) \cdot \cos(120^\circ)} \\
I_N &= \sqrt{100 + 64 +49 - 2 \cdot (80 + 56 + 70) \cdot (-0,5)} \\
I_N &= \sqrt{213 - 2 \cdot (206) \cdot (-0,5)} \\
I_N &= \sqrt{213 + 206} = \sqrt{419} = \boxed{\qty{20,47}{A}}
\end{align*}
}

\section{Representaci\'on{} gr\'afica{}}
\begin{center}
\begin{tikzpicture}[>=stealth, thick]
\draw[->] (0,0) -- (3,0) node[right] {Fase A};
\draw[->] (0,0) -- (-1.5,2.6) node[above] {Fase B};
\draw[->] (0,0) -- (-1.5,-2.6) node[below] {Fase C};
\draw[dashed, ->] (0,0) -- (0,3) node[above] {Neutro};
\end{tikzpicture}
\end{center}

\actividad{Responder el siguiente cuestionario}
\begin{enumerate}
\item \textquestiondown{}Qu\'e{} es un sistema trif\'asico{} desbalanceado?
\item \textquestiondown{}Qu\'e{} es la corriente de neutro?
\item \textquestiondown{}Por qu\'e{} en un sistema trif\'asico{} balanceado la corriente de neutro es cero?
\item \textquestiondown{}C\'omo{} se comportan las corrientes en las tres fases de un sistema trif\'asico{} balanceado?
\item \textquestiondown{}Cu\'al{} es el desfase entre las fases de un sistema trif\'asico{} ideal?
\item Explique c\'omo{} influye el desfase entre las corrientes de las fases en la corriente de neutro.
\item \textquestiondown{}Qu\'e{} consecuencias tiene el desbalance en un sistema trif\'asico{}?
\item \textquestiondown{}C\'omo{} se puede medir la corriente de neutro en un sistema desbalanceado?
\item \textquestiondown{}Qu\'e{} sucede si la corriente de neutro es demasiado elevada en un sistema el\'ectrico{}?
\item \textquestiondown{}Qu\'e{} factores pueden causar un desbalance en un sistema trif\'asico{}?
\end{enumerate}

\actividad{Resolver los siguientes problemas}
\begin{enumerate}
\item En un sistema trif\'asico{} desbalanceado, se tiene $I_A = \qty{12}{A}$, $I_B = \qty{9}{A}$, $I_C = \qty{5}{A}$. Calcula la corriente de neutro.
\item Si en un sistema trif\'asico{} desbalanceado se miden $I_A = \qty{15}{A}$, $I_B = \qty{12}{A}$, $I_C = \qty{8}{A}$, \textquestiondown{}cu\'al{} es la corriente de neutro?
\item Calcular la corriente de neutro en un sistema con $I_A = \qty{20}{A}$, $I_B = \qty{18}{A}$, $I_C = \qty{15}{A}$.
\item Un sistema trif\'asico{} tiene $I_A = \qty{25}{A}$, $I_B = \qty{20}{A}$, $I_C = \qty{17}{A}$. Determina la corriente de neutro.
\item Determina la corriente de neutro si $I_A = \qty{30}{A}$, $I_B = \qty{25}{A}$, $I_C = \qty{22}{A}$.
\item En un sistema con $I_A = \qty{8}{A}$, $I_B = \qty{6}{A}$, $I_C = \qty{4}{A}$, \textquestiondown{}cu\'al{} es la corriente de neutro?
\item Un sistema desbalanceado tiene $I_A = \qty{10}{A}$, $I_B = \qty{7}{A}$, $I_C = \qty{5}{A}$. Calcula la corriente de neutro.
\end{enumerate}

\actividad{Resolver (usar n\'umeros{} complejos para encontrar la corriente de neutro)}
\textit{Tome secuencia $ABC$ y \'angulos{} en grados. Indique $\underline{I}_A,\ \underline{I}_B,\ \underline{I}_C$ y $I_N$.}
\begin{enumerate}
\setcounter{enumi}{7}
\item $\underline{V}_{AN}=220\angle 0^\circ\ \text{V},\ \underline{V}_{BN}=228\angle -120^\circ\ \text{V},\ \underline{V}_{CN}=232\angle 120^\circ\ \text{V};\ \underline{Z}_A=18+j9\ \Omega,\ \underline{Z}_B=22+j8\ \Omega,\ \underline{Z}_C=20+j12\ \Omega$.
\item $\underline{V}_{AN}=230\angle 0^\circ\ \text{V},\ \underline{V}_{BN}=225\angle -120^\circ\ \text{V},\ \underline{V}_{CN}=220\angle 120^\circ\ \text{V};\ \underline{Z}_A=16+j12\ \Omega,\ \underline{Z}_B=24+j6\ \Omega,\ \underline{Z}_C=18+j9\ \Omega$.
\item $\underline{V}_{AN}=215\angle 0^\circ\ \text{V},\ \underline{V}_{BN}=215\angle -120^\circ\ \text{V},\ \underline{V}_{CN}=240\angle 120^\circ\ \text{V};\ \underline{Z}_A=25+j5\ \Omega,\ \underline{Z}_B=20+j10\ \Omega,\ \underline{Z}_C=15+j12\ \Omega$.
\item $\underline{V}_{AN}=225\angle 0^\circ\ \text{V},\ \underline{V}_{BN}=235\angle -118^\circ\ \text{V},\ \underline{V}_{CN}=230\angle 122^\circ\ \text{V};\ \underline{Z}_A=12+j8\ \Omega,\ \underline{Z}_B=28+j7\ \Omega,\ \underline{Z}_C=21+j9\ \Omega$.
\item $\underline{V}_{AN}=218\angle 0^\circ\ \text{V},\ \underline{V}_{BN}=222\angle -120^\circ\ \text{V},\ \underline{V}_{CN}=226\angle 120^\circ\ \text{V};\ \underline{Z}_A=30+j10\ \Omega,\ \underline{Z}_B=18+j12\ \Omega,\ \underline{Z}_C=26+j6\ \Omega$.
\item $\underline{V}_{AN}=240\angle 0^\circ\ \text{V},\ \underline{V}_{BN}=230\angle -120^\circ\ \text{V},\ \underline{V}_{CN}=235\angle 120^\circ\ \text{V};\ \underline{Z}_A=22+j11\ \Omega,\ \underline{Z}_B=16+j8\ \Omega,\ \underline{Z}_C=32+j10\ \Omega$.
\end{enumerate}
\end{document}